\documentclass[11pt,a4paper]{article}
\usepackage[margin=25mm]{geometry}
\usepackage{amsmath,amssymb,amsthm,hyperref}
\setlength{\emergencystretch}{2em}
\newtheorem{lemma}{Lemma}
\newtheorem{theorem}[lemma]{Theorem}
\newcommand{\gen}[1]{\langle#1\rangle}
\title{The central kernel in a symmetric action on pairs}
\author{Human Compression project: a component of the L044 recovery}
\date{}
\begin{document}\maketitle
The acquired faithful-lift reduction leaves the following exact case.
The natural alternating group $A_{2m}$ has a string C-marking
$\rho_0,\ldots,\rho_{m-1}$, with $m\ge8$ even. Its omission
$H=\gen{\rho_1,\ldots,\rho_{m-1}}$ is transitive on a system of
$m$ pairs and induces $S_m$. Put
\[
 B=\gen{\rho_1,\ldots,\rho_{m-2}},\quad
 V=\mathbb F_2^m,\quad z=(1,\ldots,1),
 \quad E=\{v\in V:\textstyle\sum v_i=0\}.
\]
The acquired reduction gives
\[
 B\cap V=\gen z,\qquad H\cap V=E.                 \tag{1}
\]
The images $s_1,\ldots,s_{m-2}$ of the marks of $B$ generate
$S_m$ independently. They form a string of involutions, but
no intersection property is assumed for them. The extra mark
$\rho_{m-1}$ is independent of $B$. We acquire the full actual
pair-centralizer lemma, the corrected primitive independent-string
bound, and B5's full irredundant bound: every transitive imprimitive
group of degree $q\ge9$ has maximum irredundant length at most $q-4$.
These bounds apply before any quotient IP has been established.

\begin{lemma}\label{foreststring}
Let $s_1,\ldots,s_{q-2}$ be independent involutions generating
$S_q$, where $q\ge6$. Suppose they form a string and every
maximal omission is intransitive in the natural action.
Then, up to simultaneous conjugacy and reversal, they are
\[
 (2\ 3),\quad(1\ 2)(3\ 4),\quad(4\ 5),\quad
 (5\ 6),\ldots,(q-1\ q).                           \tag{2}
\]
No intersection property is needed.
\end{lemma}
\begin{proof}
The noncommutation graph is connected. Otherwise its components
generate commuting nontrivial normal subgroups of $S_q$.
Their intersection is central, hence trivial, giving a nontrivial
direct decomposition. This is impossible: any nontrivial normal
subgroup contains $A_q$, whose centralizer in $S_q$ is trivial.
Since all far-apart marks commute, the noncommutation graph
is therefore the path on the given order of marks.

For each label choose an edge joining distinct orbits of its
omission. The acquired two-component fracture-forest lemma applies
literally: its proof uses only these orbit separations and the
matching condition on each label, not IP. In detail, the selected
$q-2$ edges form a forest, since a cycle would bypass a selected
edge with other labels. An extra edge joins the two components;
otherwise its selected-tree path bypasses some differently labelled
selected edge. Two extra edges must intersect. If they do not have
the same endpoints, the two selected paths force one alternating
square, whose labels each occur twice and which admits no other
extra edge. Thus the full graph is either a tree with one repeated
label, a tree with several repeated labels sharing one edge, or
one alternating square with uniquely labelled attached trees.
We now impose the string condition on these exact possibilities.

In the tree case with several repeated labels, write the common
edge as $\tau=\{a,b\}$. Every repeated mark is $\tau\sigma_i$,
where $\sigma_i$ is its other edge, disjoint from $\tau$.
Every edge incident to $\tau$ is a single transposition and
fails to commute with every repeated mark. Consequently there
are at most two repeated labels. If there are two, call them
$i,j$. Every connector incident to $\tau$ is a common neighbour
of $i,j$ in the noncommutation path. A path has at most one
such neighbour, so there is exactly one connector $k$. Moreover
$i,j$ commute, since otherwise $i,k,j$ would be a triangle;
hence $\sigma_i,\sigma_j$ are disjoint. Removing $k$ leaves
$\tau$ on its own two vertices, and both other edges on the
opposite side. The unique tree path from $\sigma_i$ to
$\sigma_j$ avoids $k$ and $\tau$. Its intervening edges are
single transpositions. Its consecutive labels therefore give a
noncommutation path from $i$ to $j$ avoiding $k$. Together with
$i,k,j$ this is a cycle, impossible. Multiple repeated labels
are excluded.

There is now one repeated mark, with two disjoint edges. If
there were at least two intervening tree edges between them,
those single transpositions and the repeated mark would again
form a cycle in the noncommutation graph. Thus exactly one
single connector lies between them. They form a path on four
vertices. An edge leaving a middle vertex of this path fails to
commute with both core marks, giving a triangle. An edge leaving
an outer vertex fails to commute with the double transposition.
That mark already has the connector as a neighbour, so at most
one edge leaves the core. Connectedness and $q>4$ supply exactly
one. The attached tree has no vertex of degree three, since
its three incident transpositions would give a triangle. Thus
it is a path attached at an outer vertex. This gives (2), and
the noncommutation path has only its two opposite string orders.

Finally, in the alternating-square case the two core marks
commute. Every edge leaving the square is a single transposition
noncommuting with both. There can be only one such edge $k$,
as two would give a cycle in the noncommutation graph. The
vertex $k$ already has its two core neighbours, so no further
edge can meet its outside endpoint. The connected full graph
then has only five vertices, contrary to $q\ge6$.
\end{proof}

\begin{lemma}\label{omissions}
In the standing central-kernel case, if $m\ge10$, every
maximal omission of $s_1,\ldots,s_{m-2}$ is intransitive.
\end{lemma}
\begin{proof}
Such an omission is a proper subgroup of $S_m$, with an
independent generating string of length $m-3$. If it were
transitive imprimitive, the acquired irredundant bound would
give length at most $m-4$. If it were primitive other than
$A_m,S_m$, the corrected independent-string bound would give
the same contradiction. The only possible transitive omission
is therefore $A_m$.

An interior omission is generated by two commuting nontrivial
subgroups, on the nonempty intervals on its two sides. They
are normal in the group they generate. They cannot generate
the simple group $A_m$, since simplicity would make both
that whole group, which is nonabelian.

At the first endpoint put $T=\gen{\rho_2,\ldots,\rho_{m-2}}$.
The omitted $\rho_0$ centralizes $T$. If its block image were
$A_m$, the acquired actual pair-centralizer lemma would force
$\rho_0$ to preserve the pairs: this primitive block group is
not its regular prime-dihedral exception. The whole parent
would then preserve the pairs, a contradiction.
At the other endpoint put $F=\gen{\rho_1,\ldots,\rho_{m-3}}$.
The extra mark $\rho_{m-1}$ centralizes $F$. An $A_m$ block
image has trivial centralizer in $S_m$, so that mark would
belong to $V$. A flip vector commuting with the transitive
$A_m$ is constant, hence is either $1$ or $z$. Both belong
to $B$, contrary to its independence. This excludes the
last possible transitive omission.
\end{proof}

\begin{lemma}\label{eight}
An independent connected string of six involutions generating
literal $S_8$ is conjugate to (2) or its reverse.
\end{lemma}
\begin{proof}
This is the bounded exact premise supplied with this manuscript.
Its engine is the acquired corrected primitive independent-string
engine, with the ambient group changed to literal $S_8$ and its
single target and maximum length both six. It lists every
nonidentity involution (there are 763) and extends prefixes by
involutions commuting with all but their last mark, and
noncommuting with that last mark. It rejects a prefix only if
one member is already redundant, it has already generated
$S_8$ too soon, it has reached length six, or its subgroup
index is less than $2^{6-j}$ at length $j$. Every rejection
is necessary for an independent completion of length six.
It takes conjugation orbits under the stabilizer of the entire
fixed prefix, beginning with $S_8$ itself. These operations
preserve every completion condition and therefore preserve
exhaustiveness. In particular there is no IP test, no
end-intersection test, and no centralizer-quotient pruning.

The complete visited-prefix counts are
\[
                  1,4,63,1308,934,113,3.
\]
There are exactly two ordered classes at length six, with
types $[4,6,3,3,3]$ and $[3,3,3,6,4]$. A separate exact
check verifies generation of literal $S_8$, independence
under all six omissions, all involution and string conditions,
and simultaneous conjugacy of each tuple to (2) or its
reverse. It also checks that the two ordered classes differ
and rejects a redundant repeated-generator control.
The trusted premises are the acquired enumeration argument
and GAP's exact finite group operations. The checker pins
its input bytes, rejects optimized Python, clears that variable
for children and requires completed executions with no GAP
errors. This is a complete length-six independent-string
boundary, not an inference from a C-group census.
\end{proof}

\begin{theorem}\label{main}
The standing central-kernel case is impossible, for every
even $m\ge8$.
\end{theorem}
\begin{proof}
The independent projected string has connected noncommutation
graph by the direct-decomposition argument in
Lemma~\ref{foreststring}. For $m\ge10$, Lemmas~\ref{omissions}
and~\ref{foreststring} put it in form (2); at $m=8$ use
Lemma~\ref{eight}. We retain either orientation relative to
the omitted mark $\rho_0$.

First, $\rho_0$ does not commute with $z$. Otherwise it and
$H$ would centralize $z$, so all of $A_{2m}$ would do so,
whereas its natural centralizer is trivial. Since $\rho_0$
centralizes $T=\gen{\rho_2,\ldots,\rho_{m-2}}$, we have
$z\notin T$. By (1), projection is consequently an isomorphism
from $T$ to its block image. This is actual faithfulness;
no IP is asserted for the whole projected marking.

If $\rho_0$ is next to the tail end of (2), its commuting
subgroup $T$ fixes one original pair setwise and induces
$S_{m-1}$ on the others. The acquired one-fixed-pair lemma
forces $\rho_0$ to preserve every pair, a contradiction.
We may thus place it next to the core end. Its commuting
subgroup $T$ now projects isomorphically to
\[
                 S_{\{1,2\}}\times S_{\{3,\ldots,m\}}.
                                                               \tag{3}
\]
Indeed, if $u=(1\ 2)(3\ 4)$ and $v=(4\ 5)$, then
$(uv)^3=(1\ 2)$; this recovers $(3\ 4)$ and all the
remaining path transpositions. Write $b=m-2$, so $b\ge6$
is even.

The preimage in $T$ of the pure $A_b$ factor in (3) fixes
the four points over $\{1,2\}$: its action there takes
values in the abelian flip group $C_2^2$, whereas $A_b$
is perfect. It has no fixed point over the other $b$
pairs, since it is transitive on those pairs. Its fixed
set is therefore exactly these four points. Thus
$\rho_0$ preserves the four points and their complement
separately.

On the $2b$ points the actual induced group of $T$ has
primitive block image $S_b$. The pair-centralizer lemma
forces the restriction of $\rho_0$ to preserve those
pairs. Its induced block permutation centralizes $S_b$,
so is trivial; its flip vector is constant. The restriction
is consequently either $1$ or the flip $z_b$ of all $b$
pairs. Both permutations are even because $b$ is even.
The restriction of $\rho_0$ to the four other points
must therefore also be even.

The pure $S_2$ element in (3) has an involutory preimage
in $T$, by faithfulness. On those four points it exchanges
the two original pairs, and hence is a double transposition.
The even centralizer of a double transposition in $S_4$
is the Klein four group of double transpositions. It
preserves every partition of the four points into pairs.
Since $\rho_0$ centralizes this element and is even on
these points, it preserves their two original pairs as
well. Thus all of $\rho_0$ preserves the pair system,
contradicting $\gen{H,\rho_0}=A_{2m}$. This proves the
theorem.
\end{proof}

\paragraph{Attribution and scope.}
The source supplies the invariant-module reduction in its
Section~3.2. The present argument uses the already acquired
full irredundant bound for imprimitive groups and the corrected
primitive bound to justify every projected omission. It derives
the small core directly from the string condition and the
acquired fracture forest, without quotient IP; the actual
centralizer and parity argument then closes the surviving
kernel case. The degree-eight independent-string census is
new and distinct from the earlier C-group census. No splitting
of $B$ as $C_2\times S_m$ is assumed. Together with the
accepted faithful component this recovers the full symmetric
block-image branch. General imprimitive cases, the final
strong-induction assembly and whole-source review remain
separate obligations. No whole-original, conjecture, Lean or
compression credit is granted by this component.
\end{document}
